
\subsubsection{2.1 RLHF Evaluation Methods}

RLHF follows a three-stage pipeline \citep{ouyang2022instructgpt}: supervised fine-tuning, reward model training from pairwise comparisons, and policy optimization via PPO. Evaluation focuses on preference win rates. InstructGPT achieved 85\% preference agreement versus GPT-3 base model. Constitutional AI \citep{bai2022constitutional} extended this with AI feedback for harmlessness, reporting improved safety scores while maintaining helpfulness. Both methods treat preference convergence as success---higher win rates indicate better alignment.

These benchmarks measure \textbf{unidirectional alignment} ($AI\rightarrow human)$: does the model produce preferred outputs? They do not measure \textbf{bidirectional alignment} ($human\rightarrow AI)$: do humans preserve critical evaluation capacity? High agreement could indicate quality improvement (users prefer better responses on objective tasks like math or factual QA) or homogenization (users habituate to model style on subjective tasks like creative writing or opinions where diverse preferences are legitimate). Existing metrics cannot distinguish these scenarios.

\textbf{Literature search:} We searched ACL Anthology, arXiv (cs.CL, cs.LG), and Google Scholar for ''RLHF diversity'', ''preference variance'', ''preference entropy'', ''annotator agreement entropy'', ''human feedback diversity'', and ''bidirectional alignment'' (2020-2026). We found work measuring inter-annotator agreement (Cohen's kappa) and variance in recommender systems but no prior work applying Shannon entropy to RLHF preference distributions as a bidirectional alignment metric.

WebGPT \citep{nakano2021webgpt} and OpenAI Summarization \citep{stiennon2020summarize} provide relevant precedents. OpenAI Summarization collected multi-annotator ratings on the same summary candidates, enabling variance analysis. However, published metrics focused on mean preference scores, not entropy or diversity measures. This suggests data structure for diversity measurement exists in some datasets but has not been leveraged.

\subsubsection{2.2 Bidirectional Alignment and Human Agency}

Human-AI interaction research emphasizes user agency and empowerment \citep{amershi2019guidelines}. Guidelines include ''support efficient correction'' and ''encourage granular feedback''---both requiring users maintain critical evaluation of AI outputs. However, these guidelines rely on self-reported measures (user surveys, perceived control) or interaction patterns (override rates, feedback frequency). Our entropy approach provides a behavioral proxy computable from existing preference data without new human evaluation.

Reward hacking literature \citep{skalse2022misalignment} addresses over-optimization where models exploit misspecified reward functions. Our concern is orthogonal: not model behavior pathology but user behavior homogenization. Entropy collapse would indicate alignment succeeded at making users agree but failed at preserving legitimate diversity on subjective tasks.

\subsubsection{2.3 Information-Theoretic Diversity Measures}

Shannon entropy \citep{shannon1948information} H = -$\Sigma$ p\_i log(p\_i) measures distribution uncertainty. Higher entropy indicates diversity; lower entropy indicates concentration. Entropy has been applied to ecological systems, information retrieval, and machine learning policy exploration. To our knowledge, its application to human preference distributions in RLHF evaluation has not been explored in prior work (see Section 2.1 literature search).

Prior work on preference diversity in recommender systems \citep{nguyen2014exploring} showed personalized algorithms can create filter bubbles by reducing content diversity. This parallels our concern: RLHF may reduce preference diversity by training users to converge on model-preferred responses. The key difference is measurement level: recommender systems track individual user exposure diversity; we propose population-level preference entropy as collective critical evaluation proxy.

\subsubsection{2.4 Dataset Format and Annotation Cost Trade-offs}

Pairwise comparison collection (Bradley-Terry models, Thurstone scaling) is standard in preference elicitation due to annotation efficiency: labelers compare two options (binary choice) rather than rating multiple candidates independently. Anthropic-HH \citep{bai2022constitutional} and WebGPT \citep{nakano2021webgpt} use this format: each example presents one chosen and one rejected response for a given prompt. This minimizes labeler cognitive load and enables large-scale collection (160K+ comparisons for Anthropic-HH).

However, pairwise formats create a diversity measurement incompatibility. To compute per-prompt preference entropy, we need multiple annotators rating the same response candidates---a distribution over fixed options (e.g., ''40\% prefer A, 35\% B, 25\% C'' $\rightarrow$ H $\approx$ 1.05 nats). Anthropic-HH instead provides unique response pairs per comparison (A1 vs B1, A2 vs B2), preventing aggregation into distributions.

OpenAI Summarization \citep{stiennon2020summarize} provides a counter-example: multiple annotators rated the same summary candidates (4-5 summaries per article, 3-5 annotators per summary). This multi-annotator fixed-response format enables entropy computation but incurs higher annotation cost (N annotators $\times$ M responses per prompt vs 1 annotator per pairwise comparison). Our contribution documents this format-measurement trade-off in RLHF benchmarks and proposes alternatives when multi-annotator data is unavailable.

\subsubsection{2.5 Positioning Our Work}

We introduce preference entropy as a proposed bidirectional alignment metric, distinguishing our work from:

\begin{itemize}
\item \textbf{RLHF evaluation} (InstructGPT, Constitutional AI): Measures preference agreement (unidirectional $AI\rightarrow human)$, not diversity (bidirectional human agency preservation).
\item \textbf{HCI user empowerment metrics}: Requires self-reported surveys; entropy uses behavioral data from existing benchmarks.
\item \textbf{Reward hacking detection}: Focuses on model over-optimization; we address user homogenization.
\end{itemize}

Our dataset structure documentation (pairwise-unique vs multi-annotator-fixed) is a methodological contribution: documenting data format requirements for diversity measurement in RLHF evaluation, tested on Anthropic-HH (n=100 prompts) with format-based inference for other datasets. This explains why existing benchmarks cannot retrospectively measure entropy without structural changes or alternative proxies.

